\chapter{Symmetry Diagonalises}\label{linear:ch:symmetry}

\begin{center}\itshape
Symmetry, as wide or as narrow as you may define its meaning, is one idea by which man through the ages has tried to comprehend
and create order, beauty, and perfection.
\upshape(Hermann Weyl, \emph{Symmetry}, 1952)
\end{center}

In the spring of 1896 Richard Dedekind, then sixty-four and living quietly in Brunswick, sent Georg Frobenius in Berlin a problem
he had carried around for ten years. Take a finite group $G$, attach a variable $x_g$ to each element, and form the $|G|\times|G|$
determinant $\det(x_{gh^{-1}})_{g,h\in G}$, the \emph{group determinant}. For abelian groups Dedekind knew the answer: the
determinant splits completely into linear factors, one for each ``character'' of the group. For non-abelian groups it does not,
and Dedekind had found, for the quaternion group of order $8$, irreducible factors of degree $2$ that he could not explain. He
offered the problem to Frobenius, who solved it within months. To do so he invented the character theory of finite groups, one
of the most useful tools of twentieth-century algebra, in a series of papers that started appearing the same year.

This chapter is the abelian half of that story, seen through linear algebra. The guiding principle fits in one line:
\emph{a matrix that commutes with a symmetry can be diagonalised by the symmetry}. When the symmetry group is abelian, its
characters give an eigenbasis that does not depend on the matrix at all. We start with circulant matrices, where the group is
cyclic and the eigenbasis is the discrete Fourier basis, then move to general finite abelian groups and their Cayley graphs.
Two examples carry the applications: the Paley graphs, whose eigenvalues are quadratic Gauss sums, and the Hamming cube, whose
symmetric matrices are governed by Krawtchouk polynomials. The chapter ends with Delsarte's linear programming bound, a
remarkable way of proving that codes cannot be too large by finding a positive semidefinite matrix in disguise.

The general (non-abelian) theory is Volume~2, Chapter~11; the fast algorithm for the discrete Fourier transform is Volume~4,
Chapter~7; Gauss sums appear in Volume~1, Chapter~6 and Volume~5, Chapter~11.

\begin{think}
Compute $\det\begin{pmatrix}a&b&c\\c&a&b\\b&c&a\end{pmatrix}$. Can you see, before expanding, that $a+b+c$ divides it? What
other linear factors could it have, if you allow complex coefficients?
\end{think}

\section{Circulant matrices}

Let $n\ge1$ and $\omega=e^{2\pi i/n}$. Index vectors of $\C^n$ by $\Z/n$, and let $S$ be the cyclic shift, $(Sx)_j=x_{j+1}$, so
that $S_{j,j+1}=1$ (indices modulo $n$) and $S^n=I$. A \emph{circulant} matrix is a matrix whose rows are successive cyclic shifts
of the first one:
\[
C=\begin{pmatrix}c_0&c_1&\cdots&c_{n-1}\\c_{n-1}&c_0&\cdots&c_{n-2}\\\vdots&&\ddots&\vdots\\c_1&c_2&\cdots&c_0\end{pmatrix},
\qquad C_{jk}=c_{k-j},
\]
that is, $C=\sum_{k=0}^{n-1}c_kS^k=p(S)$ with $p(z)=c_0+c_1z+\dots+c_{n-1}z^{n-1}$, the \emph{representer polynomial}.

\begin{theorem}[eigenvalues of circulants]\label{linear:thm:circulant}
For $a\in\Z/n$ let $v_a=(\omega^{aj})_{j\in\Z/n}$. Then $Cv_a=p(\omega^a)v_a$ for every circulant $C=p(S)$. The vectors
$v_a/\sqrt n$ form an orthonormal basis of $\C^n$, so every circulant is normal, and
\[
\det C=\prod_{a=0}^{n-1}p(\omega^a).
\]
\end{theorem}
\begin{proof}
$(Sv_a)_j=\omega^{a(j+1)}=\omega^a(v_a)_j$, so $Sv_a=\omega^av_a$ and $S^kv_a=\omega^{ak}v_a$; summing, $p(S)v_a=p(\omega^a)v_a$.
For $a\ne b$, $\langle v_a,v_b\rangle=\sum_j\omega^{(a-b)j}=0$, the sum of a geometric progression of ratio $\omega^{a-b}\ne1$
with $\omega^{(a-b)n}=1$; and $\|v_a\|^2=n$. So the unitary matrix $F$ with columns $v_a/\sqrt n$ satisfies $F^*CF=\operatorname
{diag}(p(\omega^a))_a$, and the determinant is the product of the diagonal.
\end{proof}

The matrix $F=\frac1{\sqrt n}(\omega^{jk})_{j,k}$ is the \emph{Fourier matrix}, and the theorem says that \emph{the discrete Fourier
transform diagonalises every circulant at once}. Figure~\ref{linear:fig:circulant} shows the eigenvalues as the values of $p$ along the
unit circle. Three consequences are immediate: circulants commute with each other; the inverse of an invertible circulant is a
circulant (it is $F\operatorname{diag}(1/p(\omega^a))F^*$, a polynomial in $S$ by Lagrange interpolation); and a circulant is
invertible exactly when $p$ has no $n$-th root of unity as a root. Products of circulants can be computed through the FFT in
$O(n\log n)$ operations, which is how computers multiply large integers and polynomials (Volume~4, Chapter~7).

\begin{figure}[ht]
\centering
\fig{linear-ch09-circulant}
\caption{The circulant with first row $(1,2,0,1,0,0,0)$ has representer $p(z)=1+2z+z^3$. Its seven eigenvalues are the values
of $p$ at the seventh roots of unity: seven points on the closed curve $\theta\mapsto p(e^{i\theta})$.}
\label{linear:fig:circulant}
\end{figure}

\begin{example}[the $3\times3$ circulant]\label{linear:exa:circ3}
With $\omega=e^{2\pi i/3}$, the determinant of the circulant with first row $(a,b,c)$ is
\[
(a+b+c)(a+b\omega+c\omega^2)(a+b\omega^2+c\omega)=a^3+b^3+c^3-3abc.
\]
The second and third factors are complex conjugate when $a,b,c$ are real, with product $a^2+b^2+c^2-ab-bc-ca=\frac12\bigl((a-b)^2+
(b-c)^2+(c-a)^2\bigr)$. So for real $a,b,c$ the sign of $a^3+b^3+c^3-3abc$ is the sign of $a+b+c$, a fact that is unpleasant to
prove by expanding. Exercise~\ref{linear:ex:circ3} finds every integer of the form $a^3+b^3+c^3-3abc$.
\end{example}

\begin{example}[cycles and their relatives]\label{linear:exa:cycle}
The cycle $C_n$ has adjacency matrix $S+S^{-1}$, with representer $z+z^{-1}$; its eigenvalues $2\cos\frac{2\pi a}n$ were
found in Example~\ref{linear:exa:spectra}. Any graph on $\Z/n$ in which $j\sim k$ depends only on $k-j$, a \emph{circulant graph},
has the eigenvalues $\sum_{s\in D}\omega^{as}$, where $D$ is the set of allowed differences. For $K_n$, $D=\Z/n\setminus\{0\}$ and
the eigenvalues are $n-1$ and, for $a\ne0$, $\sum_{s\ne0}\omega^{as}=-1$.
\end{example}

Conversely, symmetry forces the shape:

\begin{proposition}\label{linear:prop:commutant}
A matrix commutes with $S$ if and only if it is a circulant.
\end{proposition}
\begin{proof}
Circulants are polynomials in $S$. Conversely, $MS=SM$ reads $M_{j,k-1}=M_{j+1,k}$ for all $j,k$ (compare the $(j,k)$ entries:
$(MS)_{jk}=M_{j,k-1}$ and $(SM)_{jk}=M_{j+1,k}$). So $M_{jk}$ depends only on $k-j$.
\end{proof}

\section{Characters of finite abelian groups}

Circulants are the case $G=\Z/n$ of a construction that works for every finite abelian group $G$, written additively. Let $\C^G$
be the space of functions $G\to\C$ with the inner product $\langle f,h\rangle=\sum_{x\in G}f(x)\overline{h(x)}$. For $g\in G$, the
\emph{translation} $T_g$ acts by $(T_gf)(x)=f(x+g)$. The translations commute, $T_gT_h=T_{g+h}$, and each one is unitary (it
permutes the values).

\begin{definition}\label{linear:def:character}
A \emph{character} of $G$ is a group homomorphism $\chi:G\to\C^\times$, that is, $\chi(x+y)=\chi(x)\chi(y)$. The characters form
the \emph{dual group} $\widehat G$ under pointwise multiplication. The trivial character is $\mathbf 1$.
\end{definition}

If $N=|G|$, then $\chi(x)^N=\chi(Nx)=\chi(0)=1$, so the values of a character are $N$-th roots of unity, $|\chi(x)|=1$ and
$\chi^{-1}=\bar\chi$. For $G=\Z/n$, a character is determined by $\chi(1)$, an $n$-th root of unity, so the characters are the
$\chi_a(x)=\omega^{ax}$: the vectors $v_a$ of Theorem~\ref{linear:thm:circulant}.

We want to show that there are always $|G|$ characters and that they form an orthogonal basis. The quickest route avoids the
structure theorem for abelian groups and uses only linear algebra.

\begin{lemma}\label{linear:lem:commuting}
A family of pairwise commuting diagonalisable operators on a finite-dimensional complex space has a common basis of eigenvectors.
\end{lemma}
\begin{proof}
By induction on the dimension. If every operator in the family is a scalar, any basis works. Otherwise pick a non-scalar $T$
and split the space into its eigenspaces $E_\lambda$, each of smaller dimension. If $U$ commutes with $T$ and $Tx=\lambda x$,
then $TUx=UTx=\lambda Ux$, so every $E_\lambda$ is invariant under the whole family. The restriction of a diagonalisable $U$ to an
invariant subspace is diagonalisable (its minimal polynomial divides that of $U$, which has simple roots). By induction each
$E_\lambda$ has a common eigenbasis, and together they give one for the whole space.
\end{proof}

\begin{theorem}[characters of a finite abelian group]\label{linear:thm:characters}
Let $G$ be a finite abelian group.
\begin{enumerate}[label=(\roman*)]
\item For distinct characters $\chi\ne\psi$, $\langle\chi,\psi\rangle=0$; and $\langle\chi,\chi\rangle=|G|$.
\item There are exactly $|G|$ characters, and they form an orthogonal basis of $\C^G$.
\item The eigenvectors common to all translations are exactly the non-zero multiples of characters: $T_g\chi=\chi(g)\chi$.
\end{enumerate}
\end{theorem}
\begin{proof}
(i) The function $\varphi=\chi\bar\psi$ is a character, non-trivial since $\chi\ne\psi$; pick $g$ with $\varphi(g)\ne1$. Shifting
the summation variable, $\sum_x\varphi(x)=\sum_x\varphi(x+g)=\varphi(g)\sum_x\varphi(x)$, so $\sum_x\varphi(x)=0$; this is
$\langle\chi,\psi\rangle$. And $\langle\chi,\chi\rangle=\sum_x|\chi(x)|^2=|G|$.

(iii) $(T_g\chi)(x)=\chi(x+g)=\chi(g)\chi(x)$. Conversely, let $f\ne0$ satisfy $T_gf=c(g)f$ for all $g$. Then $f(g)=(T_gf)(0)=
c(g)f(0)$, so $f(0)\ne0$, and we may scale to $f(0)=1$. Then $f=c$, and $c(g+h)f=T_{g+h}f=T_gT_hf=c(g)c(h)f$ shows that $c$ is a
character.

(ii) Each $T_g$ satisfies $T_g^{|G|}=I$, so its minimal polynomial divides $z^{|G|}-1$, which has simple roots: $T_g$ is
diagonalisable. By Lemma~\ref{linear:lem:commuting} the translations have a common eigenbasis $f_1,\dots,f_{|G|}$, and by (iii) each
$f_i$ is a multiple of a character $\chi_i$. Two of them cannot be multiples of the same character, since they are linearly
independent. So there are at least $|G|$ characters. By (i) the characters are pairwise orthogonal and non-zero, hence linearly
independent in the $|G|$-dimensional space $\C^G$: there are at most $|G|$ of them.
\end{proof}

Applying (ii) to $\widehat G$ in the other direction gives the \emph{dual orthogonality}: the matrix $(\chi(x))_{x,\chi}/\sqrt{|G|}$ is
unitary, so its rows are orthonormal too, which reads $\sum_{\chi\in\widehat G}\chi(x)\overline{\chi(y)}=|G|\,[x=y]$. The
\emph{Fourier transform} of $f\in\C^G$ is $\hat f(\chi)=\langle f,\chi\rangle=\sum_xf(x)\overline{\chi(x)}$, and the expansion in the
orthogonal basis is the \emph{Fourier inversion formula}
\[
f=\frac1{|G|}\sum_{\chi\in\widehat G}\hat f(\chi)\,\chi.
\]

Now the guiding principle. Call a matrix $M=(M_{x,y})_{x,y\in G}$ \emph{translation-invariant} if $M_{x+g,y+g}=M_{x,y}$ for all
$x,y,g$, or equivalently $M_{x,y}=m(x-y)$ for a function $m$ on $G$; equivalently again, $M$ commutes with every $T_g$.

\begin{theorem}[symmetry diagonalises]\label{linear:thm:invariant}
Every translation-invariant matrix $M_{x,y}=m(x-y)$ on a finite abelian group $G$ has every character as an eigenvector:
\[
M\chi=\hat m(\chi)\,\chi,\qquad\hat m(\chi)=\sum_{u\in G}m(u)\overline{\chi(u)}.
\]
In particular $\det M=\prod_{\chi\in\widehat G}\hat m(\chi)$.
\end{theorem}
\begin{proof}
$(M\chi)(x)=\sum_ym(x-y)\chi(y)=\sum_um(u)\chi(x-u)=\chi(x)\sum_um(u)\overline{\chi(u)}$, substituting $u=x-y$ and using
$\chi(x-u)=\chi(x)\overline{\chi(u)}$. By Theorem~\ref{linear:thm:characters} the characters form a basis, so $M$ is diagonal in it.
\end{proof}

For $G=\Z/n$ this is Theorem~\ref{linear:thm:circulant} again (with $m(u)=c_{-u}$). The striking feature is that the eigenvectors do not
depend on $m$: all translation-invariant matrices are diagonalised by the same basis, and their eigenvalues are Fourier
coefficients. Products of such matrices correspond to \emph{convolutions} of the functions $m$, and the theorem becomes the
\emph{convolution theorem} $\widehat{m*m'}=\hat m\,\hat m'$ (Exercise~\ref{linear:ex:convolution}).

\begin{remark}
For a non-abelian group, the matrices commuting with all translations are no longer simultaneously diagonalisable: there are not
enough one-dimensional invariant subspaces. Frobenius's answer is to block-diagonalise instead, with one block for each
irreducible representation, repeated as many times as its dimension. Dedekind's quaternion factors of degree $2$ come from the
two-dimensional representation of the quaternion group. This is the content of Volume~2, Chapter~11.
\end{remark}

\section{Cayley graphs and the Paley graphs}

Let $D\subseteq G\setminus\{0\}$ be \emph{symmetric}, $D=-D$. The \emph{Cayley graph} $\operatorname{Cay}(G,D)$ has vertex set $G$,
with $x\sim y$ when $y-x\in D$. It is $|D|$-regular, and its adjacency matrix is translation-invariant with $m=\mathbf 1_D$. By
Theorem~\ref{linear:thm:invariant}, its eigenvalues are
\[
\lambda_\chi=\sum_{s\in D}\overline{\chi(s)}=\sum_{s\in D}\chi(s),
\]
real because $D=-D$ pairs $\chi(s)$ with $\chi(-s)=\overline{\chi(s)}$. Circulant graphs are the case $G=\Z/n$. The hypercube
$Q_n$ is $\operatorname{Cay}((\Z/2)^n,\{e_1,\dots,e_n\})$, and its characters are $\chi_T(x)=(-1)^{\sum_{i\in T}x_i}$; we come back to
it in the next section.

The prettiest example uses squares. Let $p$ be a prime with $p\equiv1\pmod4$, and let $D\subseteq\Z/p$ be the set of non-zero
squares. It has $\frac{p-1}2$ elements (each non-zero square is $x^2$ for exactly two values $\pm x$), and it is symmetric because
$-1$ is a square modulo $p$ when $p\equiv1\pmod4$ (Volume~1, Chapter~6). The \emph{Paley graph} $P_p$ is $\operatorname{Cay}(\Z/p,D)$:
two residues are friends when their difference is a square. Figure~\ref{linear:fig:paley} shows $P_{13}$.

\begin{figure}[ht]
\centering
\fig{linear-ch09-paley}
\caption{The Paley graph $P_{13}$: $x\sim y$ when $x-y$ is one of the squares $\pm1,\pm3,\pm4$ modulo $13$ (difference $\pm1$ blue,
$\pm3$ red, $\pm4$ green). It is strongly regular with parameters $(13,6,2,3)$ and eigenvalues $6$ and $\frac{-1\pm\sqrt{13}}2$.}
\label{linear:fig:paley}
\end{figure}

\begin{theorem}[Paley graphs]\label{linear:thm:paley}
Let $p\equiv1\pmod4$ be prime and $\zeta=e^{2\pi i/p}$. The eigenvalues of $P_p$ are $\frac{p-1}2$ (once) and $\frac{-1\pm\sqrt p}2$,
each with multiplicity $\frac{p-1}2$. Consequently $P_p$ is strongly regular with parameters $\bigl(p,\frac{p-1}2,\frac{p-5}4,
\frac{p-1}4\bigr)$.
\end{theorem}
\begin{proof}
The character $\chi_a(x)=\zeta^{ax}$ has eigenvalue $\lambda_a=\sum_{s\in D}\zeta^{as}$. For $a=0$ this is $|D|=\frac{p-1}2$. For
$a\ne0$, introduce the \emph{quadratic Gauss sum}
\[
g(a)=\sum_{x\in\Z/p}\zeta^{ax^2}.
\]
The term $x=0$ contributes $1$, and every $s\in D$ appears twice as $x^2$, so $g(a)=1+2\lambda_a$ and $\lambda_a=\frac{g(a)-1}2$.
We compute $|g(a)|^2$ by setting $x=y+t$:
\[
|g(a)|^2=\sum_{x,y}\zeta^{a(x^2-y^2)}=\sum_t\zeta^{at^2}\sum_y\zeta^{2aty}=p,
\]
because the inner sum is $p$ for $t=0$ and $0$ otherwise ($2at$ is invertible modulo the odd prime $p$, so $y\mapsto2aty$ runs
through $\Z/p$). Next, $g(a)$ is real: $\overline{g(a)}=\sum_x\zeta^{-ax^2}=\sum_x\zeta^{a(cx)^2}=g(a)$, where $c^2=-1$ and
$x\mapsto cx$ is a bijection. So $g(a)=\pm\sqrt p$ and $\lambda_a=\frac{-1\pm\sqrt p}2$. Finally, the trace of the adjacency matrix
is $0$: $\frac{p-1}2+\sum_{a\ne0}\frac{g(a)-1}2=0$ gives $\sum_{a\ne0}g(a)=0$, so the signs $+$ and $-$ occur equally often.

For strong regularity, let $\theta,\tau=\frac{-1\pm\sqrt p}2$ and $k=\frac{p-1}2$, so $\theta+\tau=-1$ and $\theta\tau=-\frac{p-1}4$.
The matrix $(A-\theta I)(A-\tau I)=A^2+A-\frac{p-1}4I$ vanishes on every $\chi_a$, $a\ne0$, so it is a multiple $cJ$ of the
all-ones matrix (it kills $\mathbf 1^\perp$ and is symmetric). Applying both sides to $\mathbf 1$ gives $cp=k^2+k-\frac{p-1}4=
\frac{(p-1)p}4$, so $c=\frac{p-1}4$ and
\[
A^2=\tfrac{p-1}4J-A+\tfrac{p-1}4I.
\]
Reading off the entries: two adjacent vertices have $\frac{p-1}4-1=\frac{p-5}4$ common neighbours, two non-adjacent ones have
$\frac{p-1}4$, and the diagonal is $\frac{p-1}2=k$.
\end{proof}

For $p=5$, $P_5$ is the pentagon, with eigenvalues $2$ and $\frac{-1\pm\sqrt5}2$: the golden ratio, once more. For $p=13$ we get
$\mathrm{srg}(13,6,2,3)$. Which sign goes with which $a$? Exercise~\ref{linear:ex:gausssign} shows that $g(a)=\bigl(\frac ap\bigr)g(1)$,
so the answer depends only on the sign of $g(1)$. Gauss conjectured in 1801 that $g(1)=+\sqrt p$, and it took him four more years
to prove it; in a letter to Olbers in 1805 he wrote that the solution finally came ``as lightning strikes''.

Hoffman's ratio bound (Theorem~\ref{linear:thm:hoffman}) now gives a statement about squares that looks hard to reach directly:
\[
\alpha(P_p)\le p\,\frac{-\tau}{k-\tau}=p\,\frac{(1+\sqrt p)/2}{(p+\sqrt p)/2}=\sqrt p.
\]
\emph{A set of residues in which no difference is a square has at most $\sqrt p$ elements.} Multiplying by a non-square maps
$P_p$ onto its complement (Exercise~\ref{linear:ex:paleyself}), so the same bound holds for sets in which every difference is a square.

\begin{remark}
The bound $\sqrt p$ is believed to be far from the truth: random-like behaviour of the squares suggests that the largest such sets
have only about $\log p$ elements, while the best lower bounds are of that order. Improving $\sqrt p$ at all was open for decades;
in 2021 Brandon Hanson and Giorgis Petridis proved a bound of about $\sqrt{p/2}$ by a polynomial method. The problem returns in
Volume~15.
\end{remark}

\section{The Hamming cube and Krawtchouk polynomials}

Now let $G=(\Z/2)^n$, the words of length $n$ over $\{0,1\}$ with addition modulo $2$, and write $|x|$ for the \emph{weight} of $x$
(its number of ones) and $d(x,y)=|x+y|$ for the \emph{Hamming distance}. The characters are the $2^n$ functions
\[
\chi_T(x)=(-1)^{\sum_{i\in T}x_i},\qquad T\subseteq\{1,\dots,n\}.
\]
They are characters, there are $2^n$ of them, so by Theorem~\ref{linear:thm:characters} there are no others. The matrix $H_n=(\chi_T(x))
_{x,T}$ of size $2^n$ is the \emph{Walsh--Hadamard matrix}: its entries are $\pm1$ and $H_n^{\mathsf T}H_n=2^nI$. It is the
Kronecker power $H_1^{\otimes n}$ of $H_1=\bigl(\begin{smallmatrix}1&1\\1&-1\end{smallmatrix}\bigr)$, and like the Fourier matrix it
has a fast algorithm, in $n2^n$ additions (Exercise~\ref{linear:ex:walsh}).

The cube has more symmetry than translations: permuting coordinates preserves distances. The matrices invariant under both are
the functions of the distance alone. For $0\le k\le n$, let $D_k$ be the \emph{distance-$k$ matrix}, $(D_k)_{x,y}=1$ if $d(x,y)=k$
and $0$ otherwise. So $D_0=I$, $D_1$ is the adjacency matrix of $Q_n$, and $\sum_kD_k=J$.

\begin{definition}\label{linear:def:krawtchouk}
The \emph{Krawtchouk polynomials} for length $n$ are
\[
K_k(t)=\sum_{j=0}^k(-1)^j\binom tj\binom{n-t}{k-j},\qquad k=0,1,\dots,n.
\]
\end{definition}

So $K_0=1$, $K_1(t)=n-2t$, $K_2(t)=\binom{n-t}2-t(n-t)+\binom t2$. Each $K_k$ is a polynomial of degree $k$ in $t$, with leading
coefficient $(-2)^k/k!$, as the generating function $\sum_kK_k(t)z^k=(1-z)^t(1+z)^{n-t}$ shows (Exercise~\ref{linear:ex:krawgen}).

\begin{theorem}[distances in the cube]\label{linear:thm:krawtchouk}
\begin{enumerate}[label=(\roman*)]
\item $D_k\chi_T=K_k(|T|)\,\chi_T$ for every $T$. So the eigenvalues of $D_k$ are the numbers $K_k(t)$, $0\le t\le n$, with
multiplicities $\binom nt$.
\item For every word $z$, the characters of weight $k$ sum to $\sum_{|T|=k}\chi_T(z)=K_k(|z|)$.
\item (Orthogonality) For all $k,l$,
\[
\sum_{t=0}^n\binom ntK_k(t)K_l(t)=2^n\binom nk[k=l].
\]
\end{enumerate}
\end{theorem}
\begin{proof}
(i) By Theorem~\ref{linear:thm:invariant} with $m$ the indicator of the words of weight $k$ (and real characters), the eigenvalue on
$\chi_T$ is $\sum_{|u|=k}\chi_T(u)=\sum_{|u|=k}(-1)^{|u\cap T|}$, identifying $u$ with its set of ones. A word of weight $k$ with
$j$ ones inside $T$ is chosen in $\binom{|T|}j\binom{n-|T|}{k-j}$ ways, and contributes $(-1)^j$: the sum is $K_k(|T|)$.

(ii) The same count, with the roles exchanged: $\chi_T(z)=(-1)^{|T\cap z|}$, and choosing $T$ of size $k$ with $j$ elements inside
the support of $z$ gives the same sum.

(iii) Compute $\operatorname{tr}(D_kD_l)$ in two ways. Directly, it is $\sum_{x,y}(D_k)_{x,y}(D_l)_{y,x}$, the number of pairs at
distance $k$ and at distance $l$: $2^n\binom nk$ if $k=l$, else $0$. Through the common eigenbasis of (i), it is
$\sum_TK_k(|T|)K_l(|T|)=\sum_t\binom ntK_k(t)K_l(t)$.
\end{proof}

So the $n+1$ matrices $D_0,\dots,D_n$ span a commutative algebra of symmetric matrices (the \emph{Bose--Mesner algebra} of the
Hamming scheme), all diagonalised by the Walsh--Hadamard basis, and the Krawtchouk polynomials are its table of eigenvalues. We
are ready for the application.

\section{Delsarte's linear programming bound}

A \emph{binary code} of length $n$ is a set $C\subseteq\{0,1\}^n$; its \emph{minimum distance} is the smallest distance between two
distinct codewords. A code with minimum distance $d$ corrects $\lfloor\frac{d-1}2\rfloor$ errors: balls of that radius around
codewords are disjoint. The central question of coding theory asks for $A(n,d)$, the largest size of a code of length $n$ and
minimum distance $d$. Constructions give lower bounds; upper bounds need an argument that works for \emph{every} code, including
codes with no structure at all.

In his 1973 thesis at the Philips research laboratory in Brussels, Philippe Delsarte found such an argument in the linear
algebra of the previous section. Its starting point is the \emph{distance distribution} of a code,
\[
A_i=\frac1{|C|}\#\{(x,y)\in C^2:d(x,y)=i\},\qquad i=0,\dots,n,
\]
the average number of codewords at distance $i$ from a codeword. Obviously $A_0=1$, $A_i\ge0$, $\sum_iA_i=|C|$, and $A_i=0$ for
$1\le i<d$. Delsarte saw that the $A_i$ satisfy many more, hidden, inequalities.

\begin{theorem}[Delsarte, 1973]\label{linear:thm:delsarte}
For every code $C\subseteq\{0,1\}^n$ and every $k=0,1,\dots,n$,
\[
\sum_{i=0}^nA_i\,K_k(i)\ge0.
\]
\end{theorem}
\begin{proof}
By Theorem~\ref{linear:thm:krawtchouk}(ii) and $\chi_T(x+y)=\chi_T(x)\chi_T(y)$,
\[
\sum_iA_iK_k(i)=\frac1{|C|}\sum_{x,y\in C}K_k(|x+y|)=\frac1{|C|}\sum_{|T|=k}\sum_{x,y\in C}\chi_T(x)\chi_T(y)
=\frac1{|C|}\sum_{|T|=k}\Bigl(\sum_{x\in C}\chi_T(x)\Bigr)^2\ge0.
\]
\end{proof}

In matrix language: $K_k(d(x,y))$ is the matrix $E_k=\sum_{|T|=k}\chi_T\chi_T^{\mathsf T}$, a positive semidefinite matrix (a multiple of
the orthogonal projection onto the span of the characters of weight $k$), and the inequality says $\mathbf 1_C^{\mathsf T}E_k\mathbf
1_C\ge0$. Whatever the code, its indicator vector cannot have a negative ``energy'' in any eigenspace.

\begin{corollary}[the linear programming bound]\label{linear:cor:lp}
$A(n,d)$ is at most the maximum of $1+\sum_{i=d}^na_i$ over all real $a_d,\dots,a_n\ge0$ satisfying
\[
\binom nk+\sum_{i=d}^na_iK_k(i)\ge0\qquad(k=1,\dots,n).
\]
\end{corollary}
\begin{proof}
For a code of minimum distance $d$, the numbers $a_i=A_i$ are feasible ($K_k(0)=\binom nk$), and $|C|=\sum_iA_i=1+\sum_{i\ge d}A_i$.
\end{proof}

This is a linear program, which a computer solves in an instant for $n$ in the hundreds.  But even a single
inequality can be strong. Here is the case $k=1$.

\begin{corollary}[Plotkin bound]\label{linear:cor:plotkin}
If $2d>n$, every code of length $n$ and minimum distance $d$ has $|C|\le\frac{2d}{2d-n}$.
\end{corollary}
\begin{proof}
With $K_1(i)=n-2i$, Delsarte's inequality for $k=1$ reads $n+\sum_{i\ge d}A_i(n-2i)\ge0$. For $i\ge d$, $n-2i\le n-2d<0$, so
$n\ge(2d-n)\sum_{i\ge d}A_i=(2d-n)(|C|-1)$. Hence $|C|\le1+\frac n{2d-n}=\frac{2d}{2d-n}$.
\end{proof}

\begin{example}[a perfect fit]\label{linear:exa:simplex}
For $n=7$, $d=4$, Plotkin gives $|C|\le8$. The eight words
\[
0000000,\ 1110100,\ 0111010,\ 0011101,\ 1001110,\ 0100111,\ 1010011,\ 1101001
\]
(the zero word and the seven cyclic shifts of $1110100$) are pairwise at distance exactly $4$, so $A(7,4)=8$. Equality in the proof
forces all non-zero distances to be equal to $d$ (Exercise~\ref{linear:ex:plotkinsharp}), which is the case here.
\end{example}

\begin{remark}
Delsarte's method is a principle, not a trick: find a family of positive semidefinite matrices that are functions of the
distance, and test the code's indicator vector against them. Asymptotically, the best known upper bound on the rate of binary
codes is still the one McEliece, Rodemich, Rumsey and Welch extracted in 1977 from Delsarte's linear program, by choosing clever
polynomials built from Krawtchouk polynomials. The same idea, with Bessel functions in place of Krawtchouk polynomials, gives the
Cohn--Elkies bound for sphere packings in $\R^n$, and in 2016 Maryna Viazovska found the magic function that makes it sharp in
dimension $8$, proving that the $E_8$ lattice packing is optimal; this earned her a Fields Medal in 2022. These bounds connect linear algebra with information theory.
\end{remark}

\section{A gem: Dedekind's group determinant}

We can now settle the abelian half of the problem that started the chapter.

\begin{theorem}[Dedekind]\label{linear:thm:groupdet}
Let $G$ be a finite abelian group and $(x_g)_{g\in G}$ independent variables. Then
\[
\det\bigl(x_{g-h}\bigr)_{g,h\in G}=\prod_{\chi\in\widehat G}\Bigl(\sum_{g\in G}\chi(g)\,x_g\Bigr).
\]
\end{theorem}
\begin{proof}
The matrix $M_{g,h}=x_{g-h}$ is translation-invariant with $m(u)=x_u$. By Theorem~\ref{linear:thm:invariant} its determinant is
$\prod_\chi\sum_ux_u\overline{\chi(u)}$, an identity between polynomials since it holds for all complex values of the
variables. As $\chi$ runs through $\widehat G$, so does $\bar\chi$, which gives the stated form.
\end{proof}

For $G=\Z/3$ this is Example~\ref{linear:exa:circ3}; for $G=(\Z/2)^2$ it says that
\[
\det\begin{pmatrix}a&b&c&d\\b&a&d&c\\c&d&a&b\\d&c&b&a\end{pmatrix}=(a+b+c+d)(a+b-c-d)(a-b+c-d)(a-b-c+d).
\]
Each linear factor is the eigenvalue on one character, and the determinant ``knows'' the whole dual group: the coefficients of
its factors are the character table. Frobenius's discovery was that for a non-abelian group the determinant still factors, into
irreducible polynomials whose degrees are the dimensions of the irreducible representations, each raised to the power of its own
degree.

\begin{beauty}
One basis serves every symmetric problem at once. The discrete Fourier basis diagonalises every circulant, the Walsh--Hadamard
basis every function of the Hamming distance, and the eigenvalues come out as sums of roots of unity: Gauss sums for the Paley
graphs, Krawtchouk polynomials for the cube. Delsarte turned this bookkeeping into a weapon: a code is a vector, the vector has
non-negative energy in every eigenspace, and that alone bounds the size of every code that could ever be built.
\end{beauty}

\section{Exercises}

\begin{exercise}[$\star$]\label{linear:ex:circalg}
Show that the $n\times n$ circulants form a commutative algebra of dimension $n$, closed under inverses and adjoints, and
compute $(I-tS)^{-1}$ explicitly for $t^n\ne1$.
\end{exercise}

\begin{exercise}[$\star\star$]\label{linear:ex:circ3}
Show that an integer $N$ can be written $N=a^3+b^3+c^3-3abc$ with $a,b,c\in\Z$ if and only if $N\not\equiv\pm3\pmod9$.
\end{exercise}

\begin{exercise}[$\star$]\label{linear:ex:fourier}
Show that the Fourier matrix $F$ satisfies $F^2=R$, where $R$ is the reversal $x_j\mapsto x_{-j}$, and $F^4=I$. Deduce that the
eigenvalues of $F$ are among $1,-1,i,-i$.
\end{exercise}

\begin{exercise}[$\star$]\label{linear:ex:dualgroup}
Show that $\widehat{G\times H}\cong\widehat G\times\widehat H$, with $(\chi,\psi)\mapsto\bigl((x,y)\mapsto\chi(x)\psi(y)\bigr)$, and that
$\widehat{\Z/n}\cong\Z/n$.
\end{exercise}

\begin{exercise}[$\star\star$]\label{linear:ex:convolution}
Define $(f*h)(x)=\sum_yf(y)h(x-y)$. Prove the convolution theorem $\widehat{f*h}=\hat f\,\hat h$ and Plancherel's formula
$\sum_x|f(x)|^2=\frac1{|G|}\sum_\chi|\hat f(\chi)|^2$. Show that the translation-invariant matrix with symbol $m$ acts on $f$ by
$f\mapsto m*f$.
\end{exercise}

\begin{exercise}[$\star\star$]\label{linear:ex:walsh}
Show that $H_n=H_1\otimes H_{n-1}$ (in the block sense $\bigl(\begin{smallmatrix}H_{n-1}&H_{n-1}\\H_{n-1}&-H_{n-1}\end{smallmatrix}\bigr)$
for a suitable order of the words), and deduce an algorithm computing $H_nf$ with $n2^n$ additions and subtractions.
\end{exercise}

\begin{exercise}[$\star$]\label{linear:ex:krawgen}
Show that $\sum_{k=0}^nK_k(t)z^k=(1-z)^t(1+z)^{n-t}$ for integers $0\le t\le n$. Deduce that $K_k(n-t)=(-1)^kK_k(t)$ and that
$\sum_kK_k(t)=2^n[t=0]$.
\end{exercise}

\begin{exercise}[$\star\star$]\label{linear:ex:gausssign}
Show that $g(a)=\bigl(\frac ap\bigr)g(1)$ for $a\ne0$, where $\bigl(\frac ap\bigr)$ is the Legendre symbol. Deduce that in $P_p$ the
eigenvalue $\frac{-1+\sqrt p}2$ belongs to the characters $\chi_a$ with $a$ a square if $g(1)>0$.
\end{exercise}

\begin{exercise}[$\star$]\label{linear:ex:paleyself}
Let $p\equiv1\pmod4$ and $c$ a non-square modulo $p$. Show that $x\mapsto cx$ is an isomorphism from $P_p$ to its complement.
Deduce again that $P_p$ and its complement have the same spectrum.
\end{exercise}

\begin{exercise}[$\star\star$]\label{linear:ex:paleyramsey}
Show that the clique number of $P_{17}$ is $3$, by hand or with the help of the bound $\sqrt{17}$. Deduce that the Ramsey number
$R(4,4)$ is at least $18$.
\end{exercise}

\begin{exercise}[$\star\star$]\label{linear:ex:plotkinsharp}
Show that a code achieving the Plotkin bound $|C|=\frac{2d}{2d-n}$ (with $2d>n$) has all its non-zero distances equal to $d$, and
that each coordinate takes the value $1$ in exactly half of its codewords.
\end{exercise}

\begin{exercise}[$\star\star\star$]\label{linear:ex:lpthree}
Plotkin's bound says nothing for $n=8$, $d=4$. Use Delsarte's inequalities for $k=1$ and $k=2$ to show that $A(8,4)\le16$, and find
a code of that size.
\end{exercise}

\section{Hints}

\begin{hint}{linear:ex:circalg}
Circulants are the polynomials in $S$ of degree less than $n$. For the inverse, expand $\sum_jt^jS^j$ and use $S^n=I$.
\end{hint}
\begin{hint}{linear:ex:circ3}
Try $(a,b,c)=(k+1,k,k)$, $(k+1,k+1,k)$, $(k+1,k,k-1)$. For the converse, reduce modulo $3$: $a^3+b^3+c^3-3abc\equiv(a+b+c)^3$.
\end{hint}
\begin{hint}{linear:ex:fourier}
$(F^2)_{jk}=\frac1n\sum_l\omega^{l(j+k)}$.
\end{hint}
\begin{hint}{linear:ex:dualgroup}
A character of $G\times H$ restricts to characters of $G\times\{0\}$ and $\{0\}\times H$.
\end{hint}
\begin{hint}{linear:ex:convolution}
Expand $\widehat{f*h}(\chi)$ and substitute $u=x-y$. For Plancherel, apply the inversion formula.
\end{hint}
\begin{hint}{linear:ex:walsh}
Split $x=(x_1,x')$ and $T$ according to whether $1\in T$.
\end{hint}
\begin{hint}{linear:ex:krawgen}
Expand $(1-z)^t(1+z)^{n-t}$ with the binomial theorem twice.
\end{hint}
\begin{hint}{linear:ex:gausssign}
Substitute $x\mapsto bx$ when $a=b^2c$.
\end{hint}
\begin{hint}{linear:ex:paleyself}
$cx-cy$ is a square if and only if $x-y$ is not.
\end{hint}
\begin{hint}{linear:ex:paleyramsey}
Use the symmetry: you may assume a clique contains $0$ and $1$.
\end{hint}
\begin{hint}{linear:ex:plotkinsharp}
Look at where the inequalities in the proof of Corollary~\ref{linear:cor:plotkin} can be equalities. For the coordinates, write
$\sum_{x,y\in C}d(x,y)$ coordinate by coordinate.
\end{hint}
\begin{hint}{linear:ex:lpthree}
The unknowns are $a_4,\dots,a_8$, and $K_2(i)=-4,-2,4,14,28$ for $i=4,\dots,8$. Add a quarter of the $k=2$ inequality to the
$k=1$ inequality. For the code, think of affine functions on $(\Z/2)^3$.
\end{hint}

\section{Solutions}

\begin{solution}{linear:ex:circalg}
\idea Proposition~\ref{linear:prop:commutant} turns ``circulant'' into ``polynomial in $S$'', and then everything is algebra in
$\C[S]\cong\C[z]/(z^n-1)$.
\writeup The circulants are the matrices $p(S)$ with $\deg p<n$; since $S^n=I$ they are closed under products, and they commute.
The $n$ matrices $I,S,\dots,S^{n-1}$ are linearly independent (their supports are disjoint), so the dimension is $n$. The adjoint of
$S$ is $S^{-1}=S^{n-1}$, so adjoints of circulants are circulants. If $C$ is invertible, its inverse commutes with $S$ (because $C$
does), hence is a circulant by Proposition~\ref{linear:prop:commutant}. Finally, $(I-tS)\sum_{j=0}^{n-1}t^jS^j=I-t^nS^n=(1-t^n)I$, so for
$t^n\ne1$,
\[
(I-tS)^{-1}=\frac1{1-t^n}\sum_{j=0}^{n-1}t^jS^j,
\]
a circulant with first row $\frac{(1,t,t^2,\dots,t^{n-1})}{1-t^n}$.
\end{solution}

\begin{solution}{linear:ex:circ3}
\idea The factorisation of Example~\ref{linear:exa:circ3} says $N=(a+b+c)\cdot Q$ with $Q=a^2+b^2+c^2-ab-bc-ca$. Choosing $a,b,c$
consecutive or nearly so makes $Q$ equal to $1$ or $3$, which produces almost every $N$; reduction modulo $3$ and $9$ shows what is
missed.
\writeup For $(k+1,k,k)$, $a+b+c=3k+1$ and $Q=1$, so $N=3k+1$. For $(k+1,k+1,k)$, $N=3k+2$. For $(k+1,k,k-1)$, $a+b+c=3k$ and $Q=
\frac12(1+1+4)=3$, so $N=9k$. Every integer is $\equiv1,2\pmod3$ or a multiple of $9$, unless it is $\equiv\pm3\pmod9$. Conversely, by
Fermat $a^3\equiv a\pmod3$, and $a+b+c\equiv a^3+b^3+c^3\equiv N\pmod3$. If $3\mid N$, then $3\mid a+b+c$, and
$Q=(a+b+c)^2-3(ab+bc+ca)$ is also divisible by $3$, so $9\mid N=(a+b+c)Q$. Hence $N\equiv\pm3\pmod9$ is impossible.
\end{solution}

\begin{solution}{linear:ex:fourier}
\idea Square the matrix and recognise a geometric sum.
\writeup $(F^2)_{jk}=\frac1n\sum_l\omega^{jl}\omega^{lk}=\frac1n\sum_l\omega^{l(j+k)}$, which is $1$ if $j+k\equiv0$ and $0$ otherwise:
$F^2=R$. Then $F^4=R^2=I$. If $Fv=\lambda v$ with $v\ne0$, then $\lambda^4v=F^4v=v$, so $\lambda^4=1$. (All four values occur as soon as
$n\ge5$; their multiplicities, roughly $n/4$ each, follow from the trace $\operatorname{tr}F=\frac1{\sqrt n}\sum_j\omega^{j^2}$,
a quadratic Gauss sum that Gauss evaluated in 1805; the exact multiplicities were worked out by McClellan and Parks in 1972.)
\end{solution}

\begin{solution}{linear:ex:dualgroup}
\idea A homomorphism out of a product is determined by its restrictions to the factors.
\writeup Given $\chi\in\widehat G$ and $\psi\in\widehat H$, $(x,y)\mapsto\chi(x)\psi(y)$ is a character of $G\times H$, and the map is a
homomorphism, injective because $\chi(x)\psi(0)=\chi(x)$ recovers $\chi$, and likewise $\psi$. Both sides have $|G||H|$ elements by
Theorem~\ref{linear:thm:characters}, so it is an isomorphism. For $\Z/n$, $\chi\mapsto\chi(1)$ is an isomorphism onto the group of $n$-th roots of
unity, which is cyclic of order $n$, generated by $\omega$.
\end{solution}

\begin{solution}{linear:ex:convolution}
\idea Characters turn translations into multiplications, so they turn convolutions (sums of translates) into products.
\writeup
\[
\widehat{f*h}(\chi)=\sum_x\sum_yf(y)h(x-y)\overline{\chi(x)}=\sum_yf(y)\overline{\chi(y)}\sum_uh(u)\overline{\chi(u)}=\hat f(\chi)\hat h(\chi),
\]
with $u=x-y$ and $\chi(x)=\chi(y)\chi(u)$. For Plancherel, $f=\frac1{|G|}\sum_\chi\hat f(\chi)\chi$ and the orthogonality of characters
give $\langle f,f\rangle=\frac1{|G|^2}\sum_\chi|\hat f(\chi)|^2\langle\chi,\chi\rangle=\frac1{|G|}\sum_\chi|\hat f(\chi)|^2$. Finally
$(Mf)(x)=\sum_ym(x-y)f(y)=(m*f)(x)$.
\end{solution}

\begin{solution}{linear:ex:walsh}
\idea The character table of a product is the Kronecker product of the tables, and a Kronecker product with a $2\times2$ matrix
is a ``butterfly'' step.
\writeup Order the words with the first coordinate most significant, and the sets $T$ likewise ($1\in T$ or not first). For
$x=(x_1,x')$ and $T=T_1\sqcup T'$ with $T_1\subseteq\{1\}$, $\chi_T(x)=(-1)^{x_1[1\in T]}\chi_{T'}(x')$, so $H_n$ is the block matrix
$\bigl(\begin{smallmatrix}H_{n-1}&H_{n-1}\\H_{n-1}&-H_{n-1}\end{smallmatrix}\bigr)$. To compute $H_n\binom uv$ (halves $u,v$ of length
$2^{n-1}$), compute $H_{n-1}u$ and $H_{n-1}v$, then return $\binom{H_{n-1}u+H_{n-1}v}{H_{n-1}u-H_{n-1}v}$: $2^n$ additions and subtractions
on top of two half-size transforms. The cost $c_n=2c_{n-1}+2^n$ with $c_0=0$ gives $c_n=n2^n$.
\end{solution}

\begin{solution}{linear:ex:krawgen}
\idea A word of weight $k$ chooses $j$ positions among $t$ and $k-j$ among $n-t$: that is a product of two binomial expansions.
\writeup Expanding,
\[
(1-z)^t(1+z)^{n-t}=\sum_j\binom tj(-z)^j\sum_l\binom{n-t}lz^l,
\]
and the coefficient of $z^k$ collects the terms with $l=k-j$: it is $K_k(t)$.
Replacing $t$ by $n-t$ gives $(1-z)^{n-t}(1+z)^t$, which is the same function of $-z$: so $K_k(n-t)=(-1)^kK_k(t)$. At $z=1$ the
generating function is $0^t2^{n-t}$, which is $2^n$ if $t=0$ and $0$ otherwise.
\end{solution}

\begin{solution}{linear:ex:gausssign}
\idea Multiplying $x$ by a constant multiplies $x^2$ by a square, so $g(a)$ depends only on the square class of $a$; there are two
classes, and the trace computation already showed the two values are opposite.
\writeup If $a=b^2a'$ with $b\ne0$, the substitution $x\mapsto bx$ (a bijection of $\Z/p$) gives $g(a')=\sum_x\zeta^{a'b^2x^2}=g(a)$. So
$g(a)=g(1)$ for squares $a$, and $g(a)=g(c)$ for all non-squares, with $c$ a fixed non-square. There are $\frac{p-1}2$ of each, and
$\sum_{a\ne0}g(a)=0$ (proof of Theorem~\ref{linear:thm:paley}) gives $g(c)=-g(1)$. So $g(a)=\bigl(\frac ap\bigr)g(1)$, and if $g(1)=+\sqrt p$
the eigenvalue $\frac{g(a)-1}2$ equals $\frac{-1+\sqrt p}2$ exactly for squares $a$.
\end{solution}

\begin{solution}{linear:ex:paleyself}
\idea Squares times a non-square are exactly the non-squares.
\writeup The non-zero residues form a cyclic group of even order $p-1$, and the squares are its subgroup of index $2$; so $c\cdot s$ is
a non-square for every non-zero square $s$, and vice versa. Hence for $x\ne y$, $cx-cy=c(x-y)$ is a square if and only if $x-y$
is not: $x\mapsto cx$ is a bijection of $\Z/p$ sending edges to non-edges and non-edges to edges. Isomorphic graphs have the same
spectrum; here this is visible directly, since the complement of $P_p$ has eigenvalues $p-1-k=k$ and $-1-\theta=\tau$, $-1-\tau=\theta$.
\end{solution}

\begin{solution}{linear:ex:paleyramsey}
\idea Translations and multiplication by squares are automorphisms, so a clique can be moved to contain $0$ and $1$; then only the
common neighbours of $0$ and $1$ matter, and there are only $\lambda=3$ of them.
\writeup The squares modulo $17$ are $\pm1,\pm2,\pm4,\pm8$. Hoffman's bound applied to the complement (Exercise~\ref{linear:ex:paleyself})
gives clique number at most $\lfloor\sqrt{17}\rfloor=4$; we show $4$ is impossible. Translations and multiplication by a non-zero
square $s$ are automorphisms of $P_{17}$ (they preserve squareness of differences). Given an edge $\{x,y\}$ in a clique, the map
$z\mapsto(z-x)/(y-x)$ is such an automorphism (as $y-x$ is a square), sending $x,y$ to $0,1$. The common neighbours of $0$ and $1$ are
the $z$ with $z$ and $z-1$ both squares: from the list, $z\in\{2,9,16\}$. Among $2,9,16$, the differences are $7,14,7$ (up to sign),
and $7$ and $14\equiv-3$ are not squares. So no two common neighbours of $0,1$ are adjacent: no clique of size $4$, while $\{0,1,2\}$ is a
triangle. By self-complementarity the independence number is also $3$. So $P_{17}$ is a $2$-colouring of the edges of $K_{17}$ (edges
and non-edges) with no monochromatic $K_4$, which means $R(4,4)>17$. (In fact $R(4,4)=18$, by Greenwood and Gleason, 1955, who used
exactly this graph.)
\end{solution}

\begin{solution}{linear:ex:plotkinsharp}
\idea Equality in a chain of inequalities forces every intermediate step to be an equality.
\writeup In the proof of Corollary~\ref{linear:cor:plotkin}, equality $|C|=\frac{2d}{2d-n}$ requires both $n+\sum_{i\ge d}A_i(n-2i)=0$ and
$\sum_{i\ge d}A_i(n-2i)=(n-2d)\sum_{i\ge d}A_i$, so $A_i=0$ for $i>d$: all non-zero distances equal $d$. The first equality says that the
Delsarte sum for $k=1$ vanishes; by the proof of Theorem~\ref{linear:thm:delsarte} it is $\frac1{|C|}\sum_{|T|=1}\bigl(\sum_{x\in C}\chi_T(x)
\bigr)^2$, so $\sum_{x\in C}(-1)^{x_i}=0$ for each coordinate $i$: each coordinate is $1$ on exactly half of the codewords.
\end{solution}

\begin{solution}{linear:ex:lpthree}
\idea A linear program is bounded by a non-negative combination of its constraints (this is LP duality); we look for weights that
make every coefficient of $a_4,\dots,a_8$ at least $1$.
\writeup For $n=8$, $K_1(i)=8-2i$ and $K_2(i)=\binom i2-i(8-i)+\binom{8-i}2$, which is $-4,-2,4,14,28$ for $i=4,5,6,7,8$. The two
inequalities read
\[
\begin{aligned}
2a_5+4a_6+6a_7+8a_8&\le8&&(k=1),\\
4a_4+2a_5&\le28+4a_6+14a_7+28a_8&&(k=2).
\end{aligned}
\]
Adding the first to a quarter of the second gives $a_4+\frac52a_5+4a_6+6a_7+8a_8\le15+a_6+\frac72a_7+7a_8$, that is,
\[
a_4+\tfrac52a_5+3a_6+\tfrac52a_7+a_8\le15.
\]
Every coefficient is at least $1$, so $\sum_{i\ge4}a_i\le15$ and $|C|\le16$. The \emph{extended Hamming code} (or Reed--Muller code
of order $1$) attains it: the $16$ sums of subsets of
\[
00001111,\quad 00110011,\quad 01010101,\quad 11111111.
\]
These are the truth tables of the affine functions on $(\Z/2)^3$; a non-constant affine function takes the value $1$ exactly four
times, so every non-zero codeword except $11111111$ has weight $4$, and the code, being closed under addition, has minimum distance
$4$. So $A(8,4)=16$, and the linear programming bound is sharp here.
\end{solution}

\section*{Sources and further reading}

\begin{itemize}
\item T.~Hawkins, \emph{The origins of the theory of group characters}, Archive for History of Exact Sciences 7 (1971), 142--170; and
\emph{The Mathematics of Frobenius in Context}, Springer, 2013. The Dedekind--Frobenius correspondence of 1896.
\item P.~J.~Davis, \emph{Circulant Matrices}, Wiley, 1979; R.~M.~Gray, \emph{Toeplitz and circulant matrices: a review}, Foundations and
Trends in Communications and Information Theory 2 (2006), 155--239.
\item J.-P.~Serre, \emph{A Course in Arithmetic}, Springer, 1973, Chapter VI (characters of finite abelian groups); A.~Terras,
\emph{Fourier Analysis on Finite Groups and Applications}, Cambridge University Press, 1999.
\item R.~E.~A.~C.~Paley, \emph{On orthogonal matrices}, Journal of Mathematics and Physics 12 (1933), 311--320; R.~E.~Greenwood and
A.~M.~Gleason, \emph{Combinatorial relations and chromatic graphs}, Canadian Journal of Mathematics 7 (1955), 1--7; B.~Hanson and
G.~Petridis, \emph{Refined estimates concerning sumsets contained in the roots of unity}, Proceedings of the London Mathematical
Society 122 (2021), 353--358.
\item K.~Ireland and M.~Rosen, \emph{A Classical Introduction to Modern Number Theory}, 2nd ed., Springer, 1990, Chapter~6 (Gauss sums,
and the sign of the quadratic Gauss sum).
\item P.~Delsarte, \emph{An algebraic approach to the association schemes of coding theory}, Philips Research Reports Supplements 10
(1973); F.~J.~MacWilliams and N.~J.~A.~Sloane, \emph{The Theory of Error-Correcting Codes}, North-Holland, 1977, Chapters~5 and~17;
R.~J.~McEliece, E.~R.~Rodemich, H.~Rumsey and L.~R.~Welch, \emph{New upper bounds on the rate of a code via the Delsarte--MacWilliams
inequalities}, IEEE Transactions on Information Theory 23 (1977), 157--166.
\item H.~Cohn and N.~Elkies, \emph{New upper bounds on sphere packings I}, Annals of Mathematics 157 (2003), 689--714; M.~S.~Viazovska,
\emph{The sphere packing problem in dimension 8}, Annals of Mathematics 185 (2017), 991--1015.
\item C.~Godsil and G.~Royle, \emph{Algebraic Graph Theory}, Springer, 2001, Chapter~10 (strongly regular graphs and Paley graphs).
\end{itemize}
